\section{Verification Strategy}

\subsection{Proof Structure}

The proof follows an iterative workflow. The abstract state and transitions
are defined first, then candidate safety properties and invariants are
formulated. Bounded model checking on a small finite instance exposes missing
side conditions early, and each counterexample is resolved by adding or
strengthening an invariant. Once no counterexamples remain at the chosen bound,
Veil's deductive invariant checker is used to verify that the full set of
safety properties and invariants is preserved by initialisation and every
action in the unbounded model.

\subsection{Relational Encoding}

A key design decision is the relational encoding of the order book. Rather
than modelling each side as a sorted list, the specification uses membership
predicates (\texttt{inBid}, \texttt{inAsk}) and per-order field functions
(\texttt{bidPx}, \texttt{bidQty}, \texttt{bidTs}, etc.). Best-bid and
best-ask are defined as ghost predicates quantifying over all active orders on
the same side. This encoding avoids the need for list-manipulation lemmas and
keeps all proof obligations within the decidable fragment handled by Veil's
SMT backend. The trade-off is that the specification is not directly
executable: it characterises legal matches rather than computing them.

\subsection{Invariant Design}

Thirteen invariants support the seven safety properties. The most important
families are: book disjointness (an order identifier cannot appear in both
books simultaneously), timestamp monotonicity and pairwise uniqueness (which
underpin the well-definedness of price-time priority), and trade-exclusion
(any order identifier that appears in a trade record is absent from both
books). These invariants were discovered incrementally through bounded model
checking and manual inspection of failed proof obligations.
